\newpage
\section{Literature Review}

\subsection{Market Efficiency and the Case for Price Trend Predictability}

The foundational premise of modern financial economics is the Efficient Market Hypothesis (EMH), which holds that asset prices at any given moment fully reflect all available information \parencite{fama1970}. In its weak form, the EMH asserts that historical price and volume data are already impounded in current prices, such that no mechanical trading rule based on past returns can generate systematic abnormal profits \parencite{fama1970, fama1991}. This directly challenges the intellectual basis of technical analysis, which presupposes that recurring price patterns carry predictive information about future returns.

Nevertheless, a substantial body of empirical evidence has accumulated that challenges the strict weak-form interpretation of efficiency. Pioneering work on price-based predictability established that simple technical trading rules applied to U.S. stock market data generate statistically significant returns \parencite{brock1992}. A systematic, algorithmic approach to chart pattern recognition, using nonparametric kernel regression to identify formations such as head-and-shoulders and double bottoms, finds that several technical indicators provide incremental information about future return distributions for U.S. stocks over the period 1962 to 1996 \parencite{lo2000}. While the economic magnitude of these findings is debated, they establish a clear empirical basis for the idea that price history is not informationally redundant.

Beyond chart patterns, the broader literature on return predictability documents a range of persistent cross-sectional phenomena that violate weak-form efficiency. Most prominently, stocks with strong recent performance continue to outperform over three- to twelve-month holding periods, a finding that is robust across markets and time periods \parencite{Jegadeesh1993}. Conversely, long-horizon return reversals suggest that markets systematically overreact to information \parencite{DeBondt1985}. Investor sentiment, itself partly driven by visual impressions of price trajectories, further predicts the cross-section of returns \parencite{Baker2006}. The comprehensive survey by \textcite{Park2007} of 95 empirical studies concludes that modern technical trading strategies consistently generate economic profits across a range of speculative markets, at least through the early 1990s, though profitability declines as markets become more informational efficient over time.

The reconciliation of these findings with market efficiency remains contested. From a neoclassical perspective, apparent predictability may reflect rational time-varying risk premia rather than exploitable mispricing \parencite{fama1991}. From a behavioral perspective, cognitive biases and limits to arbitrage create persistent, exploitable patterns that informed traders cannot fully eliminate. This thesis does not attempt to resolve this debate, but rather contributes to the empirical evidence on whether price chart patterns carry predictive information in European equity markets, a question that remains largely unaddressed in the existing literature.

\subsection{The Factor Zoo and the Machine Learning Turn in Asset Pricing}

The past three decades of empirical asset pricing have produced an extensive catalogue of cross-sectional return predictors, collectively termed the ``factor zoo'' \parencite{Gu2020}. Standard risk adjustment frameworks, beginning with the three-factor model of \textcite{famafrench1993} incorporating market, size, and value premia and extended by \textcite{Fama2015} to include profitability and investment factors, have formed the backbone of return decomposition. The momentum factor, validated by \textcite{Jegadeesh1993} and incorporated into the four-factor model of \textcite{Carhart1997}, has proven particularly durable. However, rigorous replication studies demonstrate that a large proportion of published anomalies fail to survive out-of-sample testing or adjustment for microcaps and transaction costs, raising serious concerns about data-snooping and multiple hypothesis testing \parencite{Hou2020}.

The inherent limitations of traditional linear econometric methods, including susceptibility to overfitting, multicollinearity, and an inability to capture non-linear feature interactions, have driven a methodological pivot toward machine learning. \textcite{Gu2020} conduct a comprehensive comparison of machine learning methods for predicting the cross-section of U.S. stock returns, demonstrating that regression trees and deep neural networks substantially outperform linear benchmarks out-of-sample. The key mechanism is the ability of non-linear models to automatically identify complex, higher-order interactions among predictors that linear specifications miss by construction. All machine learning methods in their study converge on the same dominant predictive signals: variations on momentum, liquidity, and volatility, precisely the types of signals that are visually salient in price charts.

This convergence is theoretically significant. It suggests that the information content of price chart patterns, long dismissed as the domain of practitioners rather than academics, may be genuinely predictive in ways that a flexible, data-driven methodology can recover without imposing predetermined functional forms.

\subsection{Image-Based Deep Learning and Price Pattern Recognition}

The application of Convolutional Neural Networks (CNNs) to financial chart images represents the natural synthesis of the technical analysis and machine learning literatures. CNNs, whose deep residual architectures enable training of networks hundreds of layers deep without degradation of the learning signal \parencite{he2016}, were originally developed for computer vision tasks such as image classification. Their ability to learn hierarchical spatial features, such as edges, textures, and composite shapes, from raw pixel inputs makes them well-suited to the recognition of geometric patterns in price charts.

The methodological breakthrough of applying CNNs to chart images rather than raw time-series data was formalised by \textcite{Jiang2023}, whose work provides the primary benchmark for this thesis. By rendering OHLCV data as standardised two-dimensional images and training CNNs to predict the sign of future returns, they demonstrate that image-derived forecasts yield out-of-sample trading strategies that substantially outperform traditional momentum and reversal signals, with Sharpe ratios reaching 7.2 on U.S. stocks. Crucially, their learned patterns exhibit context independence, as short-term patterns remain predictive at longer horizons and patterns extracted from U.S. stocks generalise to international markets. This finding directly motivates the transfer learning component of the present thesis.

The visual encoding of financial data carries economic content beyond what naive dimensionality might suggest. The vertical dimension of a candlestick bar naturally encodes realised intraday volatility, which is among the most accurate estimators available for daily price variability \parencite{Andersen1998}. Simultaneously, the horizontal structure of a chart window captures momentum, mean-reversion, and distributional shifts in a spatially coherent format that a CNN can recognise through its convolutional filters. This dual encoding of directional and distributional information makes chart images informationally richer inputs than raw return sequences of equivalent length.

Prior to the image-based approach, the literature on systematic pattern recognition relied on kernel smoothing methods to automate the identification of classical formations \parencite{lo2000}. While that approach provided statistically significant results, it was fundamentally constrained by the need to pre-specify the patterns being tested. The CNN approach, by contrast, discovers predictive patterns in a data-adaptive manner, with no prior assumptions about their shape or economic interpretation. This distinction is fundamental to the contribution of \textcite{Jiang2023} and to the methodological design of this thesis.

\subsection{Architectural Advances: Attention Mechanisms, Hybrid Models, and Transfer Learning}

The baseline CNN architecture of \textcite{Jiang2023} has prompted subsequent literature to explore richer model families. The introduction of attention mechanisms, formalised by \textcite{vaswani2017} in the Transformer architecture, enables models to weight different segments of the input dynamically, focusing computational resources on the most predictive regions of a chart or the most informative timesteps of a sequence. Applied to financial chart images, attention-augmented CNNs have demonstrated improved directional accuracy by learning to emphasise, for example, breakout candles or volume spikes relative to recent trading ranges.

Hybrid architectures that combine image-based CNN branches with sequential branches processing raw numerical features have attracted increasing attention. Long Short-Term Memory networks \parencite{hochreiter1997}, designed specifically to capture long-range dependencies in sequential data without the vanishing gradient problem that afflicts standard recurrent networks, complement CNN branches by encoding the precise quantitative structure of the time series alongside the visual information of the chart image. Graph Convolutional Networks extend this further by modelling cross-asset correlations and spillover effects as a learnable graph structure \parencite{Chen2021Graph}. The consistent finding across these architectures is that multimodal models outperform unimodal approaches, suggesting that visual and numerical representations of market data carry complementary rather than redundant information.

Transfer learning, defined as the reuse of model weights pre-trained on a source domain for prediction in a related target domain, has become a central technique in applied deep learning following the success of large pre-trained vision models such as ResNet \parencite{he2016}. In financial applications, transfer learning carries particular promise because the data scarcity that limits model capacity in individual markets can be mitigated by training on larger, more liquid markets and fine-tuning on the target market. \textcite{Jiang2023} provide initial evidence that U.S.-trained CNN models retain predictive power in international markets without retraining, which they interpret as evidence of universal price formation mechanisms. However, systematic evaluation of this claim across European markets, which differ from U.S. markets in market microstructure, investor composition, and regulatory environment, has not been conducted.

\subsection{The European Gap and the Contribution of This Thesis}

Despite the rapid growth of image-based CNN literature, the geographic scope of existing studies is highly skewed. The large-sample evidence overwhelmingly concerns U.S. equities, with a secondary cluster of studies focused on Chinese and other Asian markets. European equity markets have received comparatively little attention in the machine learning asset pricing literature more broadly \parencite{cakici2023}, and virtually no attention in the image-based CNN forecasting literature specifically.

This gap is not merely geographic. European equity markets exhibit a range of institutional characteristics, including greater regulatory intervention, higher cross-country heterogeneity, different sector composition, and partially segmented liquidity, that may cause the visual pattern structure of price charts to differ systematically from U.S. patterns. If CNN-learned patterns are universal in the sense of \textcite{Jiang2023}, they should transfer effectively to European data. If they are instead market-specific, locally trained models should outperform transferred ones. Resolving this empirically requires a study of the scale and methodological rigour applied in this thesis.

This thesis makes three contributions to this literature. First, it provides the first large-sample, out-of-sample evaluation of image-based CNN models trained directly on Euro Stoxx 50 constituent stocks, covering the period 2005 to 2026 with a strict walk-forward experimental design. Second, it tests whether CNN models pre-trained on U.S. market data retain predictive power in European equities, directly evaluating the context-independence claim of \textcite{Jiang2023}. Third, it benchmarks performance against the Fama-French factor models \parencite{famafrench1993, Fama2015}, providing an economically interpretable yardstick for the practical significance of any predictive gains. Together, these contributions address a clear and consequential gap in the empirical literature on machine learning in finance.